\chapter{Introducción}\label{cap:introduccion}

En los últimos años se ha podido observar un notable crecimiento en la popularidad de los juegos de cartas coleccionables (TCG) como pueden ser Pokémon, Magic: The Gathering y Yu-Gi-Oh! Por esto surge la necesidad de usar una herramienta centralizada que permite la construcción de mazos y gestión de inventario. Este trabajo describe el diseño e implementación de DeckBuilder, una aplicación que busca facilitar e integrar las tareas de los coleccionistas.

El proyecto consiste en el desarrollo de una API REST que permita operaciones CRUD (Create, Read, Update, Delete) sobre barajas e inventarios de cartas, implementado en C\# mediante .NET 8 y usando MongoDB como base de datos. La elección de estas tecnologías, más en concreto de MongoDB, permite gestionar datos heterogéneos con estructuras y atributos variables, lo que encaja perfectamente con nuestra propuesta multijuego.

El objetivo principal del proyecto es construir una aplicación capaz registrar cartas obtenidas a través de APIs externas dedicadas \cite{tcgdexApi,scryfallApi,ygoprodeckApi}, de forma que se garantice lo máximo posible la precisión de la información. De este modo, se busca resolver la necesidad de usar múltiples aplicaciones para gestionar las colecciones de distintas franquicias.

El proyecto se estructura en tres fases:

\begin{enumerate}
    \item Planificación y Diseño
    \item Desarrollo e implementación
    \item Documentación y Presentación
\end{enumerate}

Como resultado, se ha obtenido una aplicación funcional que destaca por su capacidad de adaptarse a distintos esquemas de datos gracias al uso de bases de datos documentales.
